\begin{longtable}{TD}
\caption{Execution-context terms.}
\label{tab:terms-execution-context}\\

\toprule
Term & Definition \\
\midrule
\endfirsthead

\toprule
Term & Definition \\
\midrule
\endhead

\bottomrule
\endfoot

\termrow{execution context}{
The environment and state under which a task attempt executes. Execution
context includes worker state, node pressure, dependency health, network
condition, disk pressure, cache state, tenant pressure, retry state, and other
runtime or domain-specific conditions.
}

\termrow{context factor}{
A feature of the execution context that may influence task behavior. Context
factors are modeled as influences, not automatically as proven causes.
}

\termrow{worker state}{
The local state of the worker during a task attempt, including local runtime
pressure, running work, memory state, garbage-collection pressure, and local
scheduling delay.
}

\termrow{worker pressure}{
Pressure on a worker or worker pool, such as busy worker slots, saturated local
queues, high utilization, or local scheduling delay.
}

\termrow{node pressure}{
Pressure on the host, virtual machine, container, or node on which the worker
runs. Examples include CPU saturation, memory pressure, disk pressure, and
network saturation.
}

\termrow{dependency health}{
The operational condition of an external dependency, such as latency, timeout
rate, error rate, saturation, quota pressure, or availability.
}

\termrow{dependency capacity}{
The available capacity of an external dependency, such as connection-pool
capacity, rate-limit budget, service quota, or downstream throughput.
}

\termrow{network condition}{
The state of the network path relevant to a task attempt, including round-trip
time, jitter, packet loss, retransmissions, bandwidth, timeouts, and
availability.
}

\termrow{network instability}{
A normalized feature representing unstable network behavior. It may be derived
from packet loss, jitter, retransmissions, timeout rate, bandwidth variation, or
recent network failures.
}

\termrow{replica state}{
The state of a replica, shard, peer, or replica set involved in a task attempt.
Replica state is relevant to replication, repair, synchronization, and
consistency tasks.
}

\termrow{replica health}{
A normalized feature representing the operational condition of a replica or
peer. Replica health is related to, but distinct from, replication lag.
}

\termrow{replication lag}{
An observation or signal describing how far a replica is behind another source
of truth. It may be measured in bytes, records, operations, versions, or time.
}

\termrow{disk pressure}{
Pressure on the storage subsystem, such as queue depth, read latency, write
latency, I/O wait, or saturation. Disk pressure is an execution-context factor,
not the same as disk I/O demand.
}

\termrow{cache state}{
The condition of a cache relevant to a task attempt, such as warm, cold, stale,
evicted, partially populated, or inconsistent.
}

\termrow{cache warmth}{
A normalized feature representing how prepared or populated a cache is for the
task attempt. Cache warmth can influence resource demand and latency.
}

\termrow{data locality}{
The proximity of required data to the worker or execution environment. Examples
include local, remote, cross-zone, cross-region, or external storage access.
}

\termrow{tenant pressure}{
Pressure or quota usage associated with a tenant, such as used concurrency,
rate-limit consumption, or tenant-level backlog.
}

\termrow{retry state}{
The state of the retry cycle for a task instance or attempt, including attempt
number, previous errors, backoff state, retry budget, and retry history.
}

\termrow{hidden system dynamics}{
Unobserved or partially observed system state that affects task behavior, such
as internal dependency queues, lock contention, cache eviction history, data
fragmentation, noisy neighbors, kernel behavior, or runtime effects. Hidden
dynamics are represented through residual error, uncertainty, and reduced
confidence when not captured by known features.
}

\end{longtable}
